\documentclass[12pt,a4paper]{article}

% 中文支持
\usepackage{xeCJK}
\setCJKmainfont{Songti SC}
\setCJKsansfont{Heiti SC}
\setCJKmonofont{Heiti SC}

% 数学包
\usepackage{amsmath,amssymb,amsthm}
\usepackage{bm}
\usepackage{mathtools}
\usepackage{booktabs}
\usepackage{algorithm}
\usepackage{algpseudocode}

% 页面设置
\usepackage{geometry}
\geometry{left=2.5cm,right=2.5cm,top=2.5cm,bottom=2.5cm}

% 定理环境
\newtheorem{theorem}{定理}
\newtheorem{lemma}{引理}
\newtheorem{proposition}{命题}
\newtheorem{assumption}{假设}
\newtheorem{remark}{注记}

% 自定义命令
\newcommand{\vect}[1]{\mathbf{#1}}
\newcommand{\mat}[1]{\mathbf{#1}}
\newcommand{\norm}[1]{\|#1\|_2}
\newcommand{\tr}{\text{tr}}
\newcommand{\Real}{\text{Re}}
\DeclareMathOperator{\rank}{rank}
\DeclareMathOperator*{\argmin}{arg\,min}
\DeclareMathOperator*{\argmax}{arg\,max}

\title{\textbf{Cell-Free ISAC 系统鲁棒波束成形优化\\数学推导}}
\author{}
\date{}

\begin{document}

\maketitle

%==============================================================================
\section{问题建模}\label{sec:formulation}
%==============================================================================

\paragraph{以目标为中心的专用感知关联。}
二元变量 $b_{mp}\in\{0,1\}$ 表示 AP $m$ 是否被分配为目标 $p$ 发射专用感知波形。令 $\mat{S}_p\succeq\mat{0}$ 表示该波形的协方差。当 $b_{mp}=1$ 时，AP $m$ 可以向 $\mat{S}_p$ 分配发射功率；当 $b_{mp}=0$ 时，对应的协方差子块必须为零。该指示变量不限制通信波束成形，通信仍保持全网协同。每个 AP 的发射硬件约束为
\begin{equation}
    \tr(\mat{E}_m\mat{R}_X)\leq P_{\max},\qquad \forall m\in\mathcal{M}.
\end{equation}
关联约束和专用感知资源约束写为
\begin{equation}
\sum_{m=1}^{M}b_{mp}=N_{\mathrm{req}},\qquad
\tr(\mat{E}_m\mat{S}_p)\leq P_{\max}b_{mp},\qquad \forall m,p.
\end{equation}
其中，关联约束给出每个目标的固定感知集群规模。若 $N_{\mathrm{req}}$ 表示上限而非固定预算，应使用 $\leq N_{\mathrm{req}}$，并显式引入前传或处理代价。

总发射协方差为 $\mat{R}_X=\sum_{k=1}^{K}\mat{W}_k+\sum_{p=1}^{P}\mat{S}_p$。为使感知接收端仅使用已知且面向目标的探测信号，目标 $p$ 的 PCRB 仅由 $\mat{S}_p$ 计算：
\begin{equation}\label{eq:clustered-fim}
 \mat{J}_p^{\mathrm{data}}(\mat{S}_p)
 =\frac{2}{\sigma_s^2}\Real\!\left\{\mat{D}_{p}^{H}\mat{S}_p\mat{D}_{p}\right\}.
\end{equation}
因此，关联变量控制能够贡献 PCRB 的专用感知资源，而通信协方差不获得 PCRB 增益。对固定的 $b_{mp}$，上述约束关于连续协方差变量为仿射约束；对固定协方差变量，它们关于 $b_{mp}$ 为线性约束。除非显式假设 UE 接收端能够消除专用感知波形，否则聚合专用感知协方差 $\sum_p\mat{S}_p$ 应计入通信 SINR 的干扰项。基于 Sum Big-M 发射关断的表述，以及将 FIM 仅视为 $\mat{R}_X$ 仿射函数的推导，对应于不同的建模假设，应由本文形式替换。

本章数学推导的原始 (P1) 问题决策变量为 per-AP 波束 $\vect{w}_{m,k} \in \mathbb{C}^{N_t}$、面向目标的感知协方差 $\mat{S}_p\in\mathbb{H}_+^{MN_t}$、AP-target 软分配权重 $b_{mp}\in[0,1]$，其为上述关联二元变量的连续松弛。具体形式为
\begin{align}
\min_{\substack{\{\vect{w}_{m,k}\}_{m \in \mathcal{M}, k \in \mathcal{K}}, \{\mat{S}_p\}_{p \in \mathcal{P}}, \\ \{b_{mp}\}_{m \in \mathcal{M}, p \in \mathcal{P}}}} \quad
& \sum_{m=1}^{M} \sum_{k=1}^{K} \|\vect{w}_{m,k}\|^2 + \sum_{p=1}^{P} \tr(\mat{S}_p) \tag{P1-objective} \\
\text{s.t.} \quad
& \min_{\|\Delta\vect{h}_k\| \leq \epsilon_h \|\hat{\vect{h}}_k\|} \frac{|(\hat{\vect{h}}_k + \Delta\vect{h}_k)^H \vect{w}_k|^2}{\sum_{j \neq k} |(\hat{\vect{h}}_k + \Delta\vect{h}_k)^H \vect{w}_j|^2 + (\hat{\vect{h}}_k + \Delta\vect{h}_k)^H \big(\textstyle\sum_p \mat{S}_p\big) (\hat{\vect{h}}_k + \Delta\vect{h}_k) + \sigma_c^2} \geq \gamma_k, \quad \forall k, \tag{P1-C1} \\
& \text{SINR}_{S,p} := \frac{\vect{g}_p^H \mat{S}_p \vect{g}_p}{\sigma_s^2} \geq \gamma_S^{\text{PoD}}, \quad \forall p \in \mathcal{P}, \tag{P1-C2} \\
& \tr\big((\mat{J}_p^{\text{data}}(\mat{S}_p))^{-1}\big) \leq \Gamma_{\text{Track}, p}, \quad \forall p \in \mathcal{P}, \tag{P1-C3} \\
& \tr(\mat{E}_m \mat{S}_p) \leq P_{\max} \, b_{mp}, \quad \forall m \in \mathcal{M}, \, p \in \mathcal{P}, \tag{P1-C4'a} \label{eq:p1c4a} \\
& \tr(\mat{E}_m \mat{R}_X) \leq P_{\max}, \quad \forall m \in \mathcal{M}, \tag{P1-C4'b} \label{eq:p1c4b} \\
& \sum_{m=1}^{M} b_{mp} = N_{\text{req}}, \quad \forall p \in \mathcal{P}^{\text{active}}, \tag{P1-C5'} \\
& \mat{S}_p \succeq \mat{0}, \quad \forall p \in \mathcal{P}, \tag{P1-C6} \\
& 0 \leq b_{mp} \leq 1, \quad \forall m \in \mathcal{M}, p \in \mathcal{P}. \tag{P1-C7'}
\end{align}
其中 $\vect{w}_k \triangleq [\vect{w}_{1,k}^T, \vect{w}_{2,k}^T, \ldots, \vect{w}_{M,k}^T]^T \in \mathbb{C}^{MN_t}$ 为所有 AP 到用户 $k$ 的协作合并波束向量，由 per-AP 变量堆叠得到；$\mat{R}_X \triangleq \sum_{k=1}^{K} \vect{w}_k \vect{w}_k^H + \sum_{p=1}^{P} \mat{S}_p$（在 §2-A 的 lifted 形式中精确定义）；$\mat{E}_m \in \mathbb{R}^{MN_t \times MN_t}$ 是从 $\mat{R}_X$ 中提取 AP $m$ 天线子块的对角选择矩阵。约束 (P1-C3) 中 $\mat{J}_p^{\text{data}}$ 是数据依赖的 Fisher 信息矩阵，对多维目标状态 $\boldsymbol{\theta}_p \in \mathbb{R}^{N_\theta}$（典型 ISAC 场景为 2D/3D 位置、或角度+距离的 $N_\theta \in \{2, 3, 4\}$）的原始形式为
\begin{equation}\label{eq:fim-original}
\mat{J}_p^{\text{data}}(\mat{S}_p) = \frac{2}{\sigma_s^2} \Real\!\left\{ \mat{D}_p^H \mat{S}_p \mat{D}_p \right\} \in \mathbb{H}_+^{N_\theta \times N_\theta},
\end{equation}
其中 $\mat{D}_p \triangleq \partial \vect{g}_p / \partial \boldsymbol{\theta}_p \in \mathbb{C}^{MN_t \times N_\theta}$ 是导向矢量对状态参数的偏导数矩阵；CRB 矩阵为 $\mat{C}_p = (\mat{J}_p^{\text{data}})^{-1}$，(P1-C3) 即 $\tr(\mat{C}_p) = \tr((\mat{J}_p^{\text{data}})^{-1}) \leq \Gamma_{\text{Track},p}$，等价于对 $N_\theta$ 维目标状态施以均方误差之和上界。\eqref{eq:fim-original} 重述 \eqref{eq:clustered-fim}，仅关于 $\mat{S}_p$ 仿射——通信协方差 $\{\mat{W}_k\}$ 不贡献 PCRB——故 $\mat{J}_p^{\text{data}}$ 在 lifted 域 (P2) 中保持 trace-of-inverse 形式，**不应**简化为 $\tr(\mat{F}_p \mat{S}_p) \geq \Gamma$（该简写仅在 $N_\theta = 1$ 的标量参数场景下成立，$N_\theta \geq 2$ 时需严格使用 \eqref{eq:fim-original}）；$N_\theta = 1$ 时 $\mat{D}_p = \nabla_{\boldsymbol{\theta}_p} \vect{g}_p \in \mathbb{C}^{MN_t \times 1}$、$\mat{J}_p^{\text{data}}$ 退化为标量 $\frac{2}{\sigma_s^2} \|\mat{D}_p^H \mat{S}_p^{1/2}\|^2$、约束退化为 $\mat{J}_p^{\text{data}} \geq 1/\Gamma$，即标量简写 $\tr(\mat{F}_p \mat{S}_p) \geq 1/\Gamma$，与 \eqref{eq:fim-original} 自洽。约束 (P1-C4'a)/(P1-C4'b) 均为线性约束：(P1-C4'a) 是关联门控——$b_{mp}=1$ 授权 AP $m$ 向目标 $p$ 的专用感知协方差分配至多 $P_{\max}$ 的功率，$b_{mp}=0$ 强制 $\mat{S}_p$ 的对应子块为零，但不关断该 AP 的通信发射；(P1-C4'b) 是标准硬件上限，约束总发射协方差 $\mat{R}_X$。原 Sum Big-M "硬关断 + 满功率释放"形式对应不同的建模假设，已由本形式替换。

问题 (P1) 是非凸的。其非凸性源自四条独立来源，分布在不同约束中。通信最差情况 SINR 约束 (P1-C1) 携带分式二次结构：分子 $|\vect{h}_k^H \vect{w}_k|^2$ 与分母 $\sum_{j \neq k} |\vect{h}_k^H \vect{w}_j|^2 + \vect{h}_k^H \big(\sum_p \mat{S}_p\big) \vect{h}_k + \sigma_c^2$ 在波束向量与感知协方差上均为二次函数，因此其比值在分母为正的流形上是非凸函数~[1]。叠加于其上，最差情况量词 $\forall \|\Delta\vect{h}_k\| \leq \epsilon_h \|\hat{\vect{h}}_k\|$ 是半无限的，由此得到的集合 $\{\vect{w} : \min_{\Delta\vect{h} \in \mathcal{B}_\epsilon} \text{SINR}_k(\Delta\vect{h}) \geq \gamma_k\}$ 对凸组合不封闭。位置 PCRB 约束 (P1-C3) 要求 $\tr((\mat{J}_p^{\text{data}})^{-1}) \leq \Gamma_{\text{Track}, p}$；虽然复合函数 $\tr(\mat{J}^{-1})$ 在 $\mat{J} \succ \mat{0}$ 上是严格凸函数，从而 (P1-C3) 本身定义凸集，但该约束涉及矩阵逆的非线性运算，需要通过 Schur 补转化为 LMI 以进入标准 SDP 求解器。AP-target 软分配权重 $b_{mp}$ 的二元性 (P1-C7') 在 $\sum_{m,p}(b_{mp} - b_{mp}^2) = 0$ 的等价重写下表现为凹惩罚 $g(b) = b - b^2$（$b \in [0,1]$ 上 $g \geq 0$）。关联门控 (P1-C4'a) 与每 AP 硬件上限 (P1-C4'b) 均为线性约束：左侧分别为线性迹项 $\tr(\mat{E}_m \mat{S}_p)$ 与 $\tr(\mat{E}_m \mat{R}_X)$（经 $\mat{E}_m$ 提取），右侧分别为 $P_{\max} \, b_{mp}$（耦合关联变量）与 $P_{\max}$（标准硬件上限）；两条约束在 $(\{\vect{w}\}, \{\mat{S}_p\}, b)$ 三类变量上保凸（box $0 \le b_{mp} \le 1$ 保凸）。PSD 锥 $\mat{S}_p \succeq \mat{0}$ 本身（(P1-C6)）、线性感知 SINR $\tr(\vect{g}_p \vect{g}_p^H \mat{S}_p) \geq \gamma_S^{\text{PoD}} \sigma_s^2$（(P1-C2)）以及 box 约束 $0 \leq b_{mp} \leq 1$（(P1-C7')）自身是凸的，但由于上述四条非凸/非线性来源，整体问题仍是非凸 MINLP。我们将在 \S\ref{sec:relaxation} 中通过凸化链逐项处理每个来源（协方差提升后出现的 rank-1 约束视为第五条非凸源，由 §3 双 DC 惩罚 SCA 联合处理；上述 $b_{mp}$ 二元性在 (P1) 域已由 box + DC 惩罚统一处理，故在 §3 中与 rank-1 共享同一 SCA 框架）。

%==============================================================================
\section{优化问题求解}\label{sec:relaxation}
%==============================================================================

本节沿 \S\ref{sec:formulation} 末识别出的五条非凸/非线性来源逐项求解。其中，SINR 分式与半无限不确定性经交叉相乘（\S\ref{sec:relaxation}-C step 1）和 S-Procedure（\S\ref{sec:relaxation}-C step 2）化为 LMI；PCRB 矩阵逆迹由 Schur 补无损转化为 LMI（\S\ref{sec:relaxation}-D）；AP-target 软分配权重 $b_{mp}$ 的二元性通过连续松弛 + DC 惩罚加入目标函数（\S\ref{sec:relaxation}-E），与 §3 中 rank-1 约束共享统一的双 DC 惩罚 SCA 框架。

\subsection*{A. 协方差提升 (Covariance Lifting) (P2)}

为后续半定松弛做准备，将 (P1) 中向量变量提升为协方差矩阵，即定义 $\mat{W}_k = \vect{w}_k \vect{w}_k^H \in \mathbb{H}_+^{MN_t}$、$\mat{R}_X \triangleq \sum_{k=1}^K \mat{W}_k + \sum_{p=1}^P \mat{S}_p$，其中面向目标的感知协方差 $\mat{S}_p \in \mathbb{H}_+^{MN_t}$ 已于 \S\ref{sec:formulation} 引入，得到 (P2)：
\begin{align}
\min_{\substack{\{\mat{W}_k\}_{k\in\mathcal{K}}, \{\mat{S}_p\}_{p\in\mathcal{P}}, \\ \{b_{mp}\}_{m\in\mathcal{M},p\in\mathcal{P}}}} \quad
& \sum_{k=1}^K \tr(\mat{W}_k) + \sum_{p=1}^P \tr(\mat{S}_p) \tag{P2} \\
\text{s.t.} \quad
& (\forall \|\Delta\vect{h}_k\| \leq \epsilon_h \|\hat{\vect{h}}_k\|: \tr(\hat{\mat{H}}_k(\Delta\vect{h}_k) \mat{W}_k) - \gamma_k \sum_{j \neq k} \tr(\hat{\mat{H}}_k(\Delta\vect{h}_k) \mat{W}_j) - \gamma_k \tr\Big(\hat{\mat{H}}_k(\Delta\vect{h}_k) \sum_p \mat{S}_p\Big) \geq \gamma_k \sigma_c^2), \quad \forall k, \tag{P2-C1} \\
& \tr(\vect{g}_p \vect{g}_p^H \mat{S}_p) \geq \gamma_S^{\text{PoD}} \sigma_s^2, \quad \forall p \in \mathcal{P}, \tag{P2-C2} \\
& \tr\big((\mat{J}_p^{\text{data}}(\mat{S}_p))^{-1}\big) \leq \Gamma_{\text{Track},p}, \quad \forall p \in \mathcal{P}, \tag{P2-C3} \label{eq:p2c3} \\
& \tr(\mat{E}_m \mat{S}_p) \leq P_{\max} \, b_{mp}, \quad \forall m \in \mathcal{M}, \, p \in \mathcal{P}, \tag{P2-C4'a} \\
& \tr(\mat{E}_m \mat{R}_X) \leq P_{\max}, \quad \forall m \in \mathcal{M}, \tag{P2-C4'b} \\
& \sum_{m=1}^{M} b_{mp} = N_{\text{req}}, \quad \forall p \in \mathcal{P}^{\text{active}}, \tag{P2-C5'} \\
& \mat{W}_k \succeq \mat{0}, \quad \forall k \in \mathcal{K}, \tag{P2-C6} \\
& \mat{S}_p \succeq \mat{0}, \quad \forall p \in \mathcal{P}, \tag{P2-C7} \\
& \rank(\mat{W}_k) = 1, \quad \forall k \in \mathcal{K}, \tag{P2-C8} \\
& 0 \leq b_{mp} \leq 1, \quad \forall m \in \mathcal{M}, p \in \mathcal{P}. \tag{P2-C9}
\end{align}
其中 $\mat{J}_p^{\text{data}}(\mat{S}_p) = \frac{2}{\sigma_s^2} \Real\!\left\{ \mat{D}_p^H \mat{S}_p \mat{D}_p \right\} \in \mathbb{H}_+^{N_\theta \times N_\theta}$ 为等效 MISO 感知观测模型下的数据依赖 Fisher 信息矩阵；$\mat{D}_p = \partial \vect{g}_p / \partial \boldsymbol{\theta}_p \in \mathbb{C}^{MN_t \times N_\theta}$ 为有效目标响应的导数矩阵，$\hat{\mat{H}}_k(\Delta\vect{h}_k) \triangleq (\hat{\vect{h}}_k + \Delta\vect{h}_k)(\hat{\vect{h}}_k + \Delta\vect{h}_k)^H$。该模型直接以专用感知协方差 $\mat{S}_p$ 描述信息增益，不假设向量化的多接收机观测架构。注意 (P2) 中每个 $\mat{S}_p$ 都是联合堆叠协方差，目标 $p$ 的专用波形仍可跨 AP 协作合成；关联门控 (P2-C4'a) 仅授权哪些 AP 可以为 $\mat{S}_p$ 分配功率，而 (P2-C2) 感知 SINR 与 (P2-C3) PCRB 都只由 $\mat{S}_p$ 计算。

(P1) 与 (P2) 严格等价：变量映射 $\vect{w}_k \leftrightarrow \mat{W}_k = \vect{w}_k \vect{w}_k^H$ 在一一对应（rank-1 约束保证可恢复），所有约束均通过 $\tr(\vect{x}\vect{x}^H \mat{W}) = \vect{x}^H \mat{W} \vect{x}$ 等恒等式等价改写。box 约束 (P2-C9) 替换了 (P1-C7')；DC 惩罚的引入（详见 \S\ref{sec:relaxation}-E）不改变 (P2) 的可行域。整体等价变换未引入任何松弛误差。

问题 (P2) 在 lifted 协方差域中重写 (P1)。协方差提升保留了 \S\ref{sec:formulation} 中识别出的六条来源中的四条：(P2-C1) 中的 SINR 分式结构（NC1）；(P2-C5')/(P2-C9) 软分配等式与 box 约束（NC3 + 软化）；半无限最差情况量词（NC2），在 (P2) 域被改写为关于 $\Delta\vect{h}_k$ 的二次形式，将由 \S\ref{sec:relaxation} Step 3 的 S-Procedure 处理；以及 (P2-C3) 中的 PCRB 矩阵逆（NC6），lifting 后保持 $\tr((\mat{J}_p^{\text{data}})^{-1}) \leq \Gamma_{\text{Track},p}$ 形式，将在 \S\ref{sec:relaxation}-D 由 Schur 补处理。原 NC5 波束-感知双线耦合在专用协方差形式下消失，因为 (P2-C4'a) 与 (P2-C4'b) 都是线性约束。Lifting 新增一项非凸性：rank-1 流形约束（NC4, (P2-C8)）要求 $\mat{W}_k = \vect{w}_k \vect{w}_k^H$。PSD 锥约束 (P2-C6) 与 (P2-C7) 自身是凸的，(P2-C2)（感知 SINR）在 lifted 域通过恒等式 $\vect{g}_p^H \mat{S}_p \vect{g}_p = \tr(\vect{g}_p \vect{g}_p^H \mat{S}_p)$ 变为关于 $\mat{S}_p$ 的仿射函数。(P2) 域中残留的非凸/非线性来源由 \S\ref{sec:relaxation} 的凸化链 + \S3 双 DC 惩罚 SCA 框架联合处理。

\subsection*{B. 半定松弛 (Semi-Definite Relaxation)}

(P2) 中 NC4 的处理由 SDR 实现：rank-1 流形 $\mathcal{F}_{\text{rank-1}}$ 本身非凸，SDR 把 (P2-C8) 松弛为仅 $\mat{W}_k \succeq \mat{0}$，即
\begin{equation}
\mat{W}_k \succeq \mat{0}, \quad \forall k, \tag{S1.1}
\end{equation}
可行域由 $\mathcal{F}_{\text{rank-1}}$ 扩为半正定锥 $\mathcal{F}_{\text{SDR}}$，SDR 解作为原问题的下界 $P_{\text{SDR}}^* \leq P_{\text{P2}}^*$。对于 per-user SINR + cell-free 协作 + 鲁棒不确定性联合的结构，多组 multicasting 的 $G_k \leq 2$~[2] 与鲁棒 MISO 的 $N_t \leq 2$~[3] 两个紧致条件均不直接适用，需用 Gaussian 随机化在 $\text{rank}(\mat{W}_k^*) > 1$ 时恢复可行 rank-1 解。具体地，生成 $L \in [100, 1000]$ 组高斯随机向量，从生成的候选解中选取满足约束且目标最优者，但该后处理方式并无紧致性能损失上界。

\subsection*{C. S-Procedure}

本节处理 NC1 中剩余的半无限鲁棒不确定性，分两步：先把原 SINR 鲁棒不等式交叉相乘化为二次型 (P2-C1 的等价改写、不再分式)，再借 S-Procedure 将对 $\Delta\vect{h}_k$ 的全称量化降为关于 S-Procedure 松弛变量 $\mu_k$ 的 LMI。

\textbf{Step 1 -- 交叉相乘 (Cross-multiplication).} 固定任一满足 $\|\Delta\vect{h}_k\| \leq \epsilon_h \|\hat{\vect{h}}_k\|$ 的最坏情况 $\Delta\vect{h}_k$，SINR 不等式在 $\sigma_c^2 > 0$ 保证分母恒正的条件下，可直接交叉相乘化为二次型
\begin{equation}
\tr(\hat{\mat{H}}_k \mat{W}_k) - \gamma_k \sum_{j \neq k} \tr(\hat{\mat{H}}_k \mat{W}_j) - \gamma_k \tr\Big(\hat{\mat{H}}_k \sum_p \mat{S}_p\Big) \geq \gamma_k \sigma_c^2, \tag{S2.1}
\end{equation}
这是严格等价而非近似，分母为正保证双向成立；鉴于 $\sigma_c^2$ 不含决策变量，无需一阶 Taylor 近似即可消去分式，从而避免 SCA 迭代误差的累积。

\textbf{Step 2 -- S-Procedure.} 把 $\Delta\vect{h}_k$ 恢复为变量、定义 $\mat{A}_k \triangleq \frac{1}{\gamma_k} \mat{W}_k - \sum_{j \neq k} \mat{W}_j - \sum_{p} \mat{S}_p$，(P2-C1) 改写为关于 $\Delta\vect{h}_k$ 的二次不等式
\[
\tilde{f}_1(\Delta\vect{h}) = (\hat{\vect{h}} + \Delta\vect{h})^H \mat{A}_k (\hat{\vect{h}} + \Delta\vect{h}) - \sigma_c^2 \geq 0, \quad \forall \|\Delta\vect{h}_k\| \leq \epsilon_h \|\hat{\vect{h}}_k\|,
\]
不确定集由 $\tilde{f}_2(\Delta\vect{h}) = \epsilon_h^2 \|\hat{\vect{h}}_k\|^2 - \|\Delta\vect{h}_k\|^2 \geq 0$ 描述。对该单二次约束在范数球上的全称量化，由 S-Procedure~[1] 给出一个有限的 LMI 等价条件：

\begin{lemma}\textbf{(S-Procedure).}\label{lem:s-procedure}
设 $\tilde{f}_m(\Delta\vect{h}) = (\Delta\vect{h})^H \mat{A}_m (\Delta\vect{h}) + 2\Real\{(\Delta\vect{h})^H \vect{b}_m\} + c_m$，$m \in \{1, 2\}$。则蕴含式 ``$\tilde{f}_1(\Delta\vect{h}) \geq 0 \Rightarrow \tilde{f}_2(\Delta\vect{h}) \geq 0$'' \emph{成立当且仅当} 存在 $\mu \geq 0$ 使得
\[
\begin{bmatrix} \mat{A}_2 + \mu \mat{A}_1 & \vect{b}_2 + \mu \vect{b}_1 \\ \vect{b}_2^H + \mu \vect{b}_1^H & c_2 + \mu c_1 \end{bmatrix} \succeq \mat{0},
\]
此为 S-Procedure 的原初形式~[1, \S4.3]，对 $\tilde{f}_2$ 描述凸集 $\{ \Delta\vect{h} : \|\Delta\vect{h}\|_2 \leq \epsilon_h \|\hat{\vect{h}}\|_2 \}$ 的范数球场景无额外前提要求。
\end{lemma}

由 Lemma~\ref{lem:s-procedure} 直接应用于 $\tilde{f}_1 \geq 0 \Rightarrow \tilde{f}_2 \geq 0$ 的蕴含式（其中 $\tilde{f}_2$ 描述范数球不确定集），(P2-C1) 的最坏情况鲁棒形式可写为以下 LMI：
\begin{equation}
\exists \mu_k \geq 0: \begin{bmatrix} \mat{A}_k + \mu_k \mat{I} & \mat{A}_k \hat{\vect{h}}_k \\[1.2ex] \hat{\vect{h}}_k^H \mat{A}_k & \hat{\vect{h}}_k^H \mat{A}_k \hat{\vect{h}}_k - \sigma_c^2 - \mu_k \epsilon_h^2 \|\hat{\vect{h}}_k\|^2 \end{bmatrix} \succeq \mat{0}. \tag{S3.1}
\end{equation}
新增的 $\mu_k \geq 0$ 是 S-Procedure 松弛变量。


\subsection*{D. PCRB 约束：Schur 补 LMI 变换}

PCRB 约束 (P1-C3) 在 lifted 域中保持为
\begin{equation}
\tr\big((\mat{J}_p^{\text{data}})^{-1}\big) \leq \Gamma_{\text{Track},p}, \quad \forall p \in \mathcal{P}, \tag{P2-C3}
\end{equation}
其中 $\mat{J}_p^{\text{data}} \succ \mat{0}$ 是 $\mat{S}_p$ 的仿射函数。矩阵求逆映射 $\mat{J} \mapsto \mat{J}^{-1}$ 在正定锥上严格矩阵凸，且迹算子 $\tr(\cdot)$ 是线性保序算子；因此复合函数 $\tr(\mat{J}^{-1})$ 是严格凸函数，约束 (P2-C3) 本身定义凸集。然而该约束含矩阵逆的非线性运算，为使其可直接被内点法 SDP 求解器处理，本节通过 Schur 补将其无损转化为 LMI。

引入辅助 Hermitian 矩阵变量 $\mat{M}_p \in \mathbb{H}^{N_\theta \times N_\theta}$（$N_\theta$ 为待估计目标参数维度），原约束可等价拆分为
\begin{subequations}\label{eq:crb_decouple}
\begin{align}
\mat{M}_p - (\mat{J}_p^{\text{data}})^{-1} \succeq \mat{0}, \quad \forall p \in \mathcal{P}, \label{eq:crb_mat_ineq} \\
\tr(\mat{M}_p) \leq \Gamma_{\text{Track},p}, \quad \forall p \in \mathcal{P}. \label{eq:crb_trace_ineq}
\end{align}
\end{subequations}
由于 $\mat{J}_p^{\text{data}} \succ \mat{0}$，根据 Schur 补引理，\eqref{eq:crb_mat_ineq} 等价于
\begin{equation}\label{eq:crb_lmi}
\begin{bmatrix} \mat{M}_p & \mat{I} \\ \mat{I} & \mat{J}_p^{\text{data}} \end{bmatrix} \succeq \mat{0}, \quad \forall p \in \mathcal{P}.
\end{equation}
因此，PCRB 约束 (P2-C3) 被完全、无近似地转化为关于 $(\{\mat{S}_p\}, \{\mat{M}_p\})$ 的线性矩阵不等式
\begin{subequations}\label{eq:crb_lmi_final}
\begin{align}
& \begin{bmatrix} \mat{M}_p & \mat{I} \\ \mat{I} & \mat{J}_p^{\text{data}} \end{bmatrix} \succeq \mat{0}, \quad \forall p \in \mathcal{P}, \label{eq:crb_lmi_final_a} \\
& \tr(\mat{M}_p) \leq \Gamma_{\text{Track},p}, \quad \forall p \in \mathcal{P}. \label{eq:crb_lmi_final_b}
\end{align}
\end{subequations}
该变换无松弛误差，避免了 SCA 外层迭代。

\subsection*{E. AP-target 软分配：连续松弛 + DC 惩罚}

本节彻底替代原文献的 path-loss 启发式 AP 选择方案。原文献将 (P2-C8) 的二元约束 $\{0,1\}$ 单独剥离出来、按大尺度衰落 $\text{PL}(d_{m,p})$ 排序后固定 top-$N_{\text{req}}$ AP 服务集 $\mathcal{M}_p$，再在固定集上求内层 SDP；该方案属 $K$-medoid 变种、本质 NP-hard，启发式外层循环无法给出全局最优性保证，与 §3 中"rank-1 约束用 DC 惩罚 SCA 严密恢复"的整体风格不一致。本节给出的严密替代方案如下：将 (P2-C9) 的 box 约束 $0 \leq b_{mp} \leq 1$ 与二元性 $b_{mp} \in \{0,1\}$ 合并为如下等价形式：
\begin{subequations}\label{eq:b-rewrite}
\begin{align}
& 0 \leq b_{mp} \leq 1, \tag{E.1a} \\
& g(b_{mp}) \triangleq b_{mp} - b_{mp}^2 = 0. \tag{E.1b}
\end{align}
\end{subequations}
对于 $b \in [0,1]$，$g(b) = b - b^2 \geq 0$ 恒成立，且 $g$ 是 $b$ 的凹函数（$g''(b) = -2 < 0$）；将 (E.1b) 的等式约束通过惩罚项加入目标函数：
\begin{equation}
\min \quad \underbrace{\sum_{k=1}^K \tr(\mat{W}_k) + \sum_{p=1}^P \tr(\mat{S}_p)}_{\text{原目标}} \; + \; \eta_b \underbrace{\sum_{m \in \mathcal{M}, p \in \mathcal{P}} g(b_{mp})}_{= \sum_{m,p} (b_{mp} - b_{mp}^2)} \quad \text{s.t. (P2-C1)} \sim \text{(P2-C9)} \setminus \{\text{(P2-C9)}\}, \tag{P2-DCP}
\end{equation}
其中 $\eta_b \geq 0$ 为惩罚因子；目标函数呈现"凸函数 $+$ 凹函数 × $\eta_b$"的 DC 结构，约束 (P2-C9) 仅保留 box 形式（box 集是凸的），但目标中 $-\eta_b b_{mp}^2$ 项在 $b_{mp}$ 上是凹函数，故 (P2-DCP) 不是凸优化问题，该非凸性将在 \S3 中与 rank-1 约束一起、由双 DC 惩罚 SCA 联合处理。$b_{mp} \in \{0,1\}$ 时 $g(b_{mp}) = 0$、惩罚项对二元解无任何代价，$b_{mp} \in (0,1)$ 时 $g(b_{mp}) > 0$、惩罚项随 $\eta_b$ 增大而增大；但有限惩罚仅压低而不必然消除该残差，精确恢复须由 \S3 的残差检验与固定分配恢复步骤确认。引入 $\eta_b$ 凹惩罚后 (P2-DCP) 与 \S3-A 中 rank-1 凹惩罚 $\rho(\mat{W}_k) = \tr(\mat{W}_k) - \lambda_{\max}(\mat{W}_k)$ 在数学结构上完全对偶——两者都是"凸函数 $-$ 凸函数"的 DC 形式，\S3 设计的 SCA 框架可同时承载两套惩罚：同一次 SCA 迭代内对 $\mat{W}_k$ 做特征分解、对 $b_{mp}$ 做一阶 Taylor 展开，得到关于 $(\{\mat{W}_k\}, \{b_{mp}\}, \{\mat{S}_p\}, \{\mu_k\}, \{\mat{M}_p\})$ 的统一凸 SDP 子问题。大尺度衰落启发式可用作 warm start，但对非凸 SCA 而言初始点可能影响收敛速度及到达的稳定点，不能据此声称最终解与初值无关。

\subsection*{F. 整理 (P3)：含双 DC 惩罚的最终形式}

将 §B 步骤 SDR（(S1.1)）、§C 步骤 S-Procedure（(S2.1) + (S3.1)）、§D Schur 补（\eqref{eq:crb_lmi_final}）、§E 连续松弛 + DC 惩罚（box (P2-C9) + (P2-DCP)）应用于 (P2)，得到最终问题 (P3)：
\begin{align}
\min_{\substack{\{\mat{W}_k\}, \{\mat{S}_p\}, \{\mu_k\}, \\ \{\mat{M}_p\}, \{b_{mp}\}}} \quad
& \sum_{k=1}^K \tr(\mat{W}_k) + \sum_{p=1}^P \tr(\mat{S}_p) + \eta_b \sum_{m, p} (b_{mp} - b_{mp}^2) \tag{P3} \\
\text{s.t.} \quad
& \begin{bmatrix} \mat{A}_k + \mu_k \mat{I} & \mat{A}_k \hat{\vect{h}}_k \\[1.5ex] \hat{\vect{h}}_k^H \mat{A}_k & \hat{\vect{h}}_k^H \mat{A}_k \hat{\vect{h}}_k - \sigma_c^2 - \mu_k \epsilon_h^2 \|\hat{\vect{h}}_k\|^2 \end{bmatrix} \succeq \mat{0}, \quad \forall k \tag{P3-C1} \label{eq:p3c1} \\
& \tr(\vect{g}_p \vect{g}_p^H \mat{S}_p) \geq \gamma_S^{\text{PoD}} \sigma_s^2, \quad \forall p \tag{P3-C2} \label{eq:p3c2} \\
& \begin{bmatrix} \mat{M}_p & \mat{I} \\ \mat{I} & \mat{J}_p^{\text{data}} \end{bmatrix} \succeq \mat{0}, \quad \forall p \tag{P3-C3} \label{eq:p3c3} \\
& \tr(\mat{M}_p) \leq \Gamma_{\text{Track},p}, \quad \forall p \tag{P3-C3b} \label{eq:p3c3b} \\
& \tr(\mat{E}_m \mat{S}_p) \leq P_{\max} \, b_{mp}, \quad \forall m, p \tag{P3-C4'a} \label{eq:p3c4a} \\
& \tr(\mat{E}_m \mat{R}_X) \leq P_{\max}, \quad \forall m \tag{P3-C4'b} \label{eq:p3c4b} \\
& \sum_{m=1}^{M} b_{mp} = N_{\text{req}}, \quad \forall p \in \mathcal{P}^{\text{active}} \tag{P3-C5'} \label{eq:p3c5} \\
& \mat{W}_k \succeq \mat{0}, \quad \forall k \tag{P3-C6} \label{eq:p3c6} \\
& \mat{S}_p \succeq \mat{0}, \quad \forall p \tag{P3-C7} \label{eq:p3c7} \\
& \mu_k \geq 0, \quad \forall k \tag{P3-C8} \label{eq:p3c8} \\
& 0 \leq b_{mp} \leq 1, \quad \forall m, p. \tag{P3-C9} \label{eq:p3c9}
\end{align}
其中 $\mat{A}_k = \frac{1}{\gamma_k} \mat{W}_k - \sum_{j \neq k} \mat{W}_j - \sum_{p} \mat{S}_p$、$\mat{J}_p^{\text{data}}$ 是 $\mat{S}_p$ 的仿射函数、$\mat{E}_m$ 为天线子块提取对角矩阵。(P3) 的可行域由 (P3-C1) LMI、(P3-C3)/(P3-C3b) LMI/线性不等式、(P3-C2)/(P3-C4'a)/(P3-C4'b)/(P3-C5') 线性等式/不等式、(P3-C6)/(P3-C7) 半正定锥、(P3-C8) 标量半空间、(P3-C9) box 约束的交集构成，每一集合均为凸集，故可行域是凸集；其目标函数 $\sum_k \tr(\mat{W}_k) + \sum_p \tr(\mat{S}_p) + \eta_b \sum_{m,p} (b_{mp} - b_{mp}^2)$ 在 $(\{\mat{W}_k\}, \{\mat{S}_p\})$ 上线性，在 $\{b_{mp}\}$ 上是"线性 $-\eta_b \cdot$ 二次"的 DC 形式，整体目标函数在 $(\{\mat{W}_k\}, \{\mat{S}_p\}, \{b_{mp}\})$ 上是 DC（非凸），故 (P3) 不是凸 SDP，其非凸性由 §3 双 DC 惩罚 SCA 框架求解。

\begin{theorem}[(P3) 的结构性质]\label{thm:p3-structure}
(P3) 的可行域为凸集；目标函数是关于 $(\{\mat{W}_k\}, \{b_{mp}\}, \{\mat{S}_p\})$ 的 DC 函数；约束矩阵关于 $(\{\mat{W}_k\}, \{\mat{S}_p\})$ 仿射、关于 $\{b_{mp}\}$ 线性；故 (P3) 是 DC 规划（difference of convex program），可由 \S3 双 DC 惩罚 SCA 框架求稳定点。
\end{theorem}

\begin{proof}
详见附录 \ref{app:proof-p3}。
\end{proof}

%==============================================================================
\section{双 DC 惩罚 SCA 框架：秩一与二元联合恢复}\label{sec:sca}
%==============================================================================

(P3) 的可行域凸、目标函数 DC；其全局最优难以直接求取，但 \emph{任何满足 (P3) Karush-Kuhn-Tucker (KKT) 条件的稳定点都是工程可接受的解}。本节设计一个双 DC 惩罚 SCA 框架：(P3) 目标中的 $-\eta_b b_{mp}^2$ 凹项与"rank-1 缺失"对应的 $-\eta_{\text{rank}} \lambda_{\max}(\mat{W}_k)$ 凹项在同一 SCA 循环内被同步线性化，每次迭代求解一个标准凸 SDP。本节同时给出 Gaussian randomization / 主特征向量提取作为基准后处理。

\subsection*{A. 双 DC 惩罚的统一表征}

对任意半正定矩阵 $\mat{W}_k \succeq \mat{0}$，其迹等于最大特征值当且仅当秩至多为 1；若 $\mat{W}_k\neq0$（例如正 SINR 目标下的 active user），这才等价于秩为 1：
\begin{equation}\label{eq:rank1-equivalence}
\tr(\mat{W}_k) - \lambda_{\max}(\mat{W}_k) = 0 \quad \Longleftrightarrow \quad \rank(\mat{W}_k) \leq 1.
\end{equation}
由于 $\lambda_{\max}(\mat{W}_k) \leq \tr(\mat{W}_k)$ 对 PSD 矩阵恒成立，定义秩亏量
\begin{equation}\label{eq:rho-def}
\rho(\mat{W}_k) \triangleq \tr(\mat{W}_k) - \lambda_{\max}(\mat{W}_k) \geq 0.
\end{equation}
对 $b \in [0,1]$，定义二元亏量
\begin{equation}\label{eq:g-def}
g(b) \triangleq b - b^2 = b(1-b) \geq 0, \quad g''(b) = -2 < 0 \;\; (\text{严格凹}).
\end{equation}
$\rho(\mat{W}_k)$ 与 $g(b_{mp})$ 在数学结构上完全同构——均为"凸函数 $-$ 凸函数"的 DC 形式，且在期望的"满足约束"集合上为零、在违反约束时为正；这是本节双 DC 惩罚能共享同一 SCA 框架的数学基础。在 (P3) 目标上额外加入 rank-1 惩罚 $\eta_{\text{rank}} \sum_k \rho(\mat{W}_k)$，得到
\begin{equation}\label{eq:p3-dcp}
\min \quad \underbrace{\sum_k \tr(\mat{W}_k) + \sum_p \tr(\mat{S}_p)}_{\text{原目标}} + \underbrace{\eta_{\text{rank}} \sum_k \rho(\mat{W}_k)}_{\text{秩一惩罚 (DC)}} + \underbrace{\eta_b \sum_{m,p} g(b_{mp})}_{\text{二元惩罚 (DC)}} \quad \text{s.t. (P3-C1)} \sim \text{(P3-C9)}.
\end{equation}
(P3-DCP) 的约束全部凸（继承自 (P3) 的凸可行域），唯一非凸性来自目标中两个 DC 惩罚项。

\subsection*{B. SCA 线性化：双 DC 的同步一阶 Taylor 展开}

在第 $t$ 次迭代，设当前迭代点为 $(\{\mat{W}_k^{(t)}\}, \{b_{mp}^{(t)}\})$。对 $\mat{W}_k^{(t)}$ 做特征分解
\begin{equation}
\mat{W}_k^{(t)} = \sum_{i=1}^{r_k} \lambda_{k,i} \, \vect{u}_{k,i} \vect{u}_{k,i}^H, \quad \lambda_{k,1} \geq \lambda_{k,2} \geq \cdots \geq \lambda_{k,r_k} > 0,
\end{equation}
其中主特征向量为 $\vect{u}_{\max,k}^{(t)} \triangleq \vect{u}_{k,1}$。由 $\lambda_{\max}(\mat{W}_k)$ 的凸性，其一阶 Taylor 展开给出全局下界：
\begin{equation}\label{eq:lambda-lb}
\lambda_{\max}(\mat{W}_k) \geq \lambda_{\max}(\mat{W}_k^{(t)}) + \tr\Big( \vect{u}_{\max,k}^{(t)} (\vect{u}_{\max,k}^{(t)})^H (\mat{W}_k - \mat{W}_k^{(t)}) \Big) = \tr\Big( \vect{u}_{\max,k}^{(t)} (\vect{u}_{\max,k}^{(t)})^H \mat{W}_k \Big).
\end{equation}
将上式代入 (P3-DCP) 的 rank-1 惩罚即得全局上界（凹项被上界替代，子问题求极小对应真实目标下界）：
\begin{equation}\label{eq:rho-linearized}
\rho(\mat{W}_k) = \tr(\mat{W}_k) - \lambda_{\max}(\mat{W}_k) \leq \tr(\mat{W}_k) - \tr\Big( \vect{u}_{\max,k}^{(t)} (\vect{u}_{\max,k}^{(t)})^H \mat{W}_k \Big) \triangleq \hat{\rho}^{(t)}(\mat{W}_k).
\end{equation}
对 $g(b) = b - b^2$ 在 $b^{(t)}$ 处作一阶 Taylor 展开（凹函数的一阶 Taylor 给全局过估计）：
\begin{equation}\label{eq:g-linearized}
g(b) \leq g(b^{(t)}) + g'(b^{(t)})(b - b^{(t)}) = b^{(t)} - (b^{(t)})^2 + (1 - 2b^{(t)})(b - b^{(t)}),
\end{equation}
去除常数项 $b^{(t)} - (b^{(t)})^2 - (1 - 2b^{(t)}) b^{(t)} = (b^{(t)})^2$（该常数不影响最小化）后，线性化惩罚为
\begin{equation}\label{eq:g-linearized-clean}
g(b) \leq (1 - 2b^{(t)}) \, b + \text{const} \triangleq \hat{g}^{(t)}(b).
\end{equation}
将 \eqref{eq:rho-linearized} 与 \eqref{eq:g-linearized-clean} 代入 (P3-DCP)，得到第 $t$ 次迭代的凸 SDP 子问题
\begin{equation}\label{eq:p3-sca}
\begin{aligned}
\min_{\substack{\{\mat{W}_k\}, \{\mat{S}_p\}, \{\mu_k\}, \\ \{\mat{M}_p\}, \{b_{mp}\}}} \quad & \sum_{k=1}^K \tr(\mat{W}_k) + \sum_{p=1}^P \tr(\mat{S}_p) + \eta_{\text{rank}} \sum_{k=1}^K \Big[ \tr(\mat{W}_k) - \tr\big( \vect{u}_{\max,k}^{(t)} (\vect{u}_{\max,k}^{(t)})^H \mat{W}_k \big) \Big] \\
& \qquad + \eta_b \sum_{m, p} (1 - 2b_{mp}^{(t)}) \, b_{mp} \\
\text{s.t.} \quad \text{(P3-C1)}\sim\text{(P3-C9)}.
\end{aligned}
\end{equation}
子问题 \eqref{eq:p3-sca} 的目标函数关于 $(\{\mat{W}_k\}, \{\mat{S}_p\}, \{b_{mp}\})$ 是线性函数（$\eta_{\text{rank}}$ 项中 $\tr(\mat{W}_k)$ 与 $-\tr(\vect{u}_{\max,k}^{(t)}(\vect{u}_{\max,k}^{(t)})^H \mat{W}_k)$ 合并为 $\tr((\mat{I} - \vect{u}_{\max,k}^{(t)}(\vect{u}_{\max,k}^{(t)})^H) \mat{W}_k)$；$\eta_b$ 项为 $b_{mp}$ 的线性函数），约束 (P3-C1)--(P3-C9) 全部为 LMI/线性/半正定锥/box，故 (P3-SCA-t) 是标准凸 SDP，可用 MOSEK/SDPT3 等内点法求解器高效求解。更新迭代点 $(\{\mat{W}_k^{(t+1)}\}, \{b_{mp}^{(t+1)}\}, \{\mat{S}_p^{(t+1)}\}, \{\mu_k^{(t+1)}\}, \{\mat{M}_p^{(t+1)}\})$ 后重复该过程。

\begin{algorithm}[htbp]
\caption{双 DC 惩罚 SCA 迭代（秩一 + 二元联合恢复）}
\label{alg:sca-joint}
\begin{algorithmic}[1]
\Require 固定惩罚因子 $\eta_{\text{rank}},\eta_b>0$、容差 $\varepsilon>0$、最大迭代次数 $T_{\max}$。
\Ensure 近似秩一协方差、满足基数约束的二元分配，以及恢复成功时经物理约束验证的解。
\State 求解无惩罚 SDR ($\eta_{\text{rank}}=\eta_b=0$) 得到可行 SCA 起点 $(\{\mat{W}_k^{(0)}\},\{b_{mp}^{(0)}\},\ldots)$；若不可行则报告 infeasible。启发式只可作 warm start，且可能影响到达的稳定点。
\For{$t = 0, 1, \ldots, T_{\max}-1$}
    \For{每个 $k \in \mathcal{K}$}
        \State 对 $\mat{W}_k^{(t)}$ 特征分解，得主特征向量 $\vect{u}_{\max,k}^{(t)}$。
    \EndFor
    \For{每个 $(m, p)$}
        \State 计算 $g'$ 系数 $c_{mp}^{(t)} = 1 - 2 b_{mp}^{(t)}$。
    \EndFor
    \State 求解凸 SDP \eqref{eq:p3-sca} 得 $(\{\mat{W}_k^{(t+1)}\}, \{b_{mp}^{(t+1)}\}, \{\mat{S}_p^{(t+1)}\}, \{\mu_k^{(t+1)}\}, \{\mat{M}_p^{(t+1)}\})$。
    \State 计算 $\delta_{\text{rank}}^{(t)} = \sum_k \rho(\mat{W}_k^{(t+1)})$、$\delta_b^{(t)} = \sum_{m,p} g(b_{mp}^{(t+1)})$。
    \State 计算真实固定惩罚目标 $F_\eta^{(t+1)}$（即 \eqref{eq:p3-dcp} 的目标）。
    \If{$\delta_{\text{rank}}^{(t)} \leq \varepsilon$、$\delta_b^{(t)} \leq \varepsilon$ 且 $|F_\eta^{(t+1)}-F_\eta^{(t)}|\leq\varepsilon$}
        \State \textbf{break}。
    \EndIf
\EndFor
\State 对每个 target $p$，将 $\{b_{mp}^{(t+1)}\}_{m=1}^M$ 中最大的 $N_{\text{req}}$ 个置为 1，其余置为 0，得到 $b_{mp}^{\rm fix}$。
\State 固定 $b=b^{\rm fix}$ 后重解连续恢复问题，并核验所有原始物理约束。若 $\delta_{\rm rank}\leq\varepsilon$ 且恢复的 $\mat{W}_k\neq0$，令 $\vect{w}_k^\star=\sqrt{\lambda_{\max}(\mat{W}_k)}\vect{u}_{\max,k}$；否则报告恢复失败。
\State 仅返回通过物理约束验证的恢复解。
\end{algorithmic}
\end{algorithm}

\begin{theorem}[双 DC SCA 收敛性]\label{thm:sca-conv}
固定 $\eta_{\text{rank}},\eta_b>0$，每次凸子问题精确求解且初始点可行时，Algorithm~\ref{alg:sca-joint} 生成的真实固定惩罚目标序列 $\{F_\eta^{(t)}\}$ 单调非增且有下界，因此收敛。本结论本身不声称任意有限惩罚会使残差严格为零。若任一聚点对所有非零 $\mat{W}_k^\star$ 满足 $\rho(\mat{W}_k^\star)=0$、对所有 $(m,p)$ 满足 $g(b_{mp}^\star)=0$，则秩一提取和保持基数的分配投影与该点一致，并满足原问题 (P1-C1)--(P1-C7')。
\end{theorem}

\begin{proof}
请见附录 \ref{app:proof-sca-conv}。
\end{proof}

\subsection*{C. 基线：Gaussian 随机化 / 主特征向量提取}

若追求单次 SDP 的极速运行，可直接求解 $\eta_{\text{rank}}=\eta_b=0$ 的 (P3) 主问题，再由主特征向量生成 rank-1 候选。这只作为基线：特征提取和分配投影可能破坏鲁棒 SINR/PCRB 保证，故每个候选必须按原问题逐项验算；不可行候选须重优化或丢弃。

\begin{algorithm}[htbp]
\caption{秩一波束提取（基线后处理）}
\label{alg:rank1}
\begin{algorithmic}[1]
\Require 收敛解 $\{\mat{W}_k^\star, \{\mat{S}_p^\star\}, b_{mp}^\star\}$。
\Ensure 秩一波束向量 $\{\vect{w}_k^\star\}$。
\For{每个 $k \in \mathcal{K}$}
    \State 计算特征分解 $\mat{W}_k^\star = \mat{U}_k \mat{\Lambda}_k \mat{U}_k^H$。
    \State 提取主导特征向量：$\tilde{\vect{w}}_k = \sqrt{\lambda_{\max}(\mat{\Lambda}_k)} \cdot \mat{u}_k^{\max}$。
    \State 缩放以满足每 AP 功率：$\vect{w}_k^\star = \eta_k \tilde{\vect{w}}_k$，$\eta_k$ 使 $\tr(\mat{E}_m \mat{R}_X) \leq P_{\max}$（其中 $\mat{R}_X = \sum_k \vect{w}_k^\star (\vect{w}_k^\star)^H + \sum_p \mat{S}_p^\star$）对所有 $m$ 成立。
\EndFor
\If{$\{\vect{w}_k^\star, b_{mp}^\star\}$ 满足所有 (P1-C1)--(P1-C7')}
    \State \Return $\{\vect{w}_k^\star, b_{mp}^\star\}$。
\Else
    \State 求解固定分配恢复问题并重新核验原始约束；若仍不可行，则丢弃该候选。
\EndIf
\end{algorithmic}
\end{algorithm}

经过 \S\ref{sec:relaxation} 的凸化链 + \S3 双 DC 惩罚 SCA，(P3) 中的所有物理约束（鲁棒 SINR、PCRB、感知 SINR、AP 功率、AP-target 关联、box）均以凸约束表示。其余秩与二元残差在同一 SCA 循环内受控；论文应报告这些数值残差、保持基数的投影和后续物理可行性核验，而不能由未经证明的精确惩罚结论直接推断恢复成功。

\appendix
\section{Proof of Theorem~\ref{thm:p3-structure}}\label{app:proof-p3}

本附录给出定理 \ref{thm:p3-structure} 的证明。

(P3) 的可行域由 (P3-C1) LMI、(P3-C2)/(P3-C3b)/(P3-C4'a)/(P3-C4'b)/(P3-C5') 线性等式/不等式、(P3-C3) LMI、(P3-C6)/(P3-C7) 半正定锥、(P3-C8) 标量半空间、(P3-C9) box 约束的交集构成。LMI 集、线性不等式/等式集、半正定锥、标量半空间、box 集均为凸集；有限个凸集的交集仍是凸集，故 (P3) 的可行域是凸集。

(P3) 的目标函数 $F(\{\mat{W}_k\}, \{b_{mp}\}, \{\mat{S}_p\}) = \sum_k \tr(\mat{W}_k) + \sum_p \tr(\mat{S}_p) + \eta_b \sum_{m,p} (b_{mp} - b_{mp}^2)$ 关于 $(\{\mat{W}_k\}, \{\mat{S}_p\})$ 线性、关于 $\{b_{mp}\}$ 是"线性 $- \eta_b \cdot$ 二次"的 DC 形式。约束矩阵关于 $(\{\mat{W}_k\}, \{\mat{S}_p\})$ 仿射、关于 $\{b_{mp}\}$ 线性，故 (P3) 是 DC 规划。

就复杂度而言，决策变量总数 $n = K \cdot (MN_t)^2 + P \cdot (MN_t)^2 + K + P \cdot N_\theta^2 + MP$（$K$ 个 $\mat{W}_k \in \mathbb{H}_+^{MN_t}$、$P$ 个 $\mat{S}_p \in \mathbb{H}_+^{MN_t}$、$K$ 个 S-Procedure 松弛变量 $\mu_k$、$P$ 个辅助矩阵 $\mat{M}_p \in \mathbb{H}_+^{N_\theta \times N_\theta}$、$MP$ 个标量 $b_{mp}$）；不等式约束数 $m = K + P + P + P + MP + M + P_{\text{active}} + K + P + K + 2MP = 3K + 4P + M + 3MP + P_{\text{active}}$（(P3-C1): $K$；(P3-C2): $P$；(P3-C3): $P$；(P3-C3b): $P$；(P3-C4'a): $MP$；(P3-C4'b): $M$；(P3-C5'): $P_{\text{active}}$；(P3-C6): $K$；(P3-C7): $P$；(P3-C8): $K$；(P3-C9): $2MP$）；最大 LMI 约束尺寸 $\ell = \max\{MN_t + 1, 2 N_\theta\}$。每次 SCA 子问题 (P3-SCA-t) 是标准凸 SDP，由内点法在多项式时间求解。

\section{Proof of Theorem~\ref{thm:sca-conv}}\label{app:proof-sca-conv}

本附录证明真实固定惩罚目标的单调性；不对任意有限惩罚下必然精确秩一/二元恢复作出断言。

令 $F_\eta(x)$ 表示 \eqref{eq:p3-dcp} 的目标，$x$ 汇集所有决策变量。第 $t$ 次迭代中，令 $Q_t(x)$ 表示由 \eqref{eq:rho-linearized} 与 \eqref{eq:g-linearized} 构造的 SCA 上界，并保留 $g$ 的完整仿射上界常数项；\eqref{eq:p3-sca} 省略该常数仅因它不改变子问题最优解。

由 $\lambda_{\max}$ 的凸性及 $g$ 的凹性，对每个可行 $x$ 都有 $F_\eta(x)\leq Q_t(x)$，且上界在当前点处紧，即 $F_\eta(x^{(t)})=Q_t(x^{(t)})$。若 $x^{(t+1)}$ 是凸子问题的最优解，则
\[
 F_\eta(x^{(t+1)})\leq Q_t(x^{(t+1)})\leq Q_t(x^{(t)})=F_\eta(x^{(t)}).
\]
故 $\{F_\eta(x^{(t)})\}$ 单调非增。功率项与两个亏量均非负，故其下界为零，从而该序列收敛。

若某一聚点满足 $\rho(\mat{W}_k^\star)=0$ 且 $\mat{W}_k^\star\neq0$，则 $\mat{W}_k^\star$ 秩一并可由主特征向量表示；若同时 $g(b_{mp}^\star)=0$，则 $b_{mp}^\star\in\{0,1\}$。只有在这一精确恢复情形下，(P3) 的凸化等价链才直接保证原问题约束。对于有限惩罚的实际计算，应报告残差，并按 Algorithm~\ref{alg:sca-joint} 做保持基数的投影、固定分配恢复和物理可行性核验。

\begin{thebibliography}{9}
\bibitem{Boyd2004}
S. Boyd and L. Vandenberghe, \emph{Convex Optimization}. Cambridge, U.K.: Cambridge University Press, 2004.

\bibitem{Luo2010}
Z.-Q. Luo, W.-K. Ma, A. M.-C. So, Y. Ye, and S. Zhang, ``Semidefinite relaxation of quadratic optimization problems,'' \emph{IEEE Signal Process. Mag.}, vol. 27, no. 3, pp. 20--34, May 2010.

\bibitem{Huang2010}
Y. Huang and D. P. Palomar, ``Rank-constrained separable semidefinite programming with applications to optimal beamforming,'' \emph{IEEE Trans. Signal Process.}, vol. 58, no. 2, pp. 664--678, Feb. 2010.
\end{thebibliography}

\end{document}
